% TEMPLATE-MILC-v3.tex - plantilla maestra UNICA de las guias MILC v3.
%
% La usan las cuatro familias de guias, cada una con su driver:
%   · Grados 6-11          -> scripts/build-guias-g11.py
%   · Bebras               -> scripts/build-guias-bebras.py
%   · Semillero            -> scripts/build-guias-semillero.py
%   · Territorio interior  -> scripts/build-guias-territorio-interior.py
%
% Antes habia tres copias casi identicas (~400 lineas iguales, 6 distintas):
% cada mejora habia que aplicarla tres veces, y en la practica no se hacia
% (el motor de recursos vivia solo en dos de ellas). Lo que varia entre
% familias esta parametrizado en 6 placeholders que cada driver llena
% (se nombran aqui SIN su sintaxis real: el builder reemplaza por texto plano,
% tambien dentro de los comentarios, y un valor multilinea rompe el preambulo):
%
%   HEADER_IZQ        encabezado izquierdo de pagina
%   HEADER_DER        encabezado derecho de pagina
%   PORTADA_ETIQUETA  rotulo sobre el numero grande (GRADO/MOMENTO/MODULO)
%   PORTADA_SERIE     linea de serie de la portada
%   PORTADA_ID        identificador de la guia en la portada
%   PORTADA_CIRCULO   el circulito de grado (°), o vacio si no aplica
%   PDF_TITULO        titulo en texto plano (sin \textbf ni comillas TeX) para
%                     los metadatos del PDF (/Title); lo produce lib_xelatex.texto_pdf
%
% Composicion (auditoria 2026-09-03): las bandas de estacion piden espacio
% (\Needspace*) y prohiben el salto justo despues (\nopagebreak), asi no quedan
% huerfanas al pie; las cajas largas (softbox, drawbox, infoband, puentebox) son
% `breakable`, asi una caja de media pagina ya no arrastra todo a la siguiente
% dejando la anterior al 40 %. Todo texto sobre mostaza o verde va en milcNegro
% (blanco sobre #E5B400 da 1,93:1 y sobre #58A923 2,95:1; ninguno pasa WCAG AA).
% El resto de claves las documenta content/guias/_SCHEMA.md. El script valida
% que no queden placeholders sin reemplazar antes de escribir el archivo final.

% Etiquetado PDF (latex-lab): /Lang, estructura y texto alternativo de las
% figuras, para lectores de pantalla. Debe ir ANTES de \documentclass.
\DocumentMetadata{lang=es-CO,tagging=on}
\documentclass[11pt,letterpaper]{article}
% headheight=24pt: la cabecera derecha lleva dos lineas (identidad de la guia
% y estacion actual). headsep se acorta en la misma medida para que el bloque
% de texto quede exactamente donde estaba con headheight=12pt.
\usepackage[margin=2.0cm,headheight=24pt,headsep=13pt]{geometry}
\usepackage[table]{xcolor}
\usepackage{fontspec}
\usepackage[spanish]{babel}
\usepackage{tikz}
% Librerías TikZ para los diagramas de las guías (bloque `recursos:`):
%  · babel  → babel en español vuelve activos caracteres como «>», y TikZ los usa
%             en sus opciones (>=stealth, ->). Sin esta librería, cualquier
%             diagrama con flechas revienta con «\language@active@arg> has an
%             extra }». Debe ir ANTES de que se dibuje nada.
%  · positioning        → sintaxis «below left=2pt and 3pt».
%  · decorations.pathreplacing → llaves ({brace}) para señalar rangos.
\usetikzlibrary{babel,positioning,decorations.pathreplacing}
\usepackage{graphicx}
% Circuitos: el trabajo de electrónica se hace hoy con micro:bit y sus sensores
% integrados (MakeCode, sin cableado), así que no se necesitan esquemáticos.
% Si algún día entran montajes con protoboard, basta descomentar esta línea:
% los diagramas .tex/.tikz declarados en `recursos:` ya se incrustan solos.
% \usepackage{circuitikz}
\usepackage[most]{tcolorbox}
\usepackage{tabularx}
\usepackage{array}
\usepackage{fancyhdr}
\usepackage{titlesec}
\usepackage[document]{ragged2e}  % bandera a la derecha en todo el cuerpo
\usepackage{enumitem}
\usepackage{microtype}
\usepackage{etoolbox}
\usepackage{needspace}
% hyperref al final del preambulo: metadatos del PDF (titulo, autor, idioma,
% familia), marcadores por estacion (\pdfbookmark) y enlaces sin recuadro.
\usepackage[hidelinks,
  pdftitle={La broma que no es broma: cuando el color y la plata tuvieron precio de lista},
  pdfauthor={PhD. Álvaro Cárdenas Orozco},
  pdfsubject={Territorio interior · Educación socioemocional · Grado 9° · Momento 5 · MILC v3},
  pdflang={es-CO},
  bookmarksopen=true,bookmarksnumbered=false]{hyperref}

% Helvetica Neue no trae la flecha U+2192, asi que cada «→» escrito literalmente
% en el YAML se compilaba a NADA, en silencio (462 flechas en 61 guias: rutas del
% tipo «paleta → tipografia → lockup» salian rotas). Mapeamos el caracter a su
% version matematica, que si tiene glifo. Asi el autor puede seguir escribiendo
% «→» comodamente en el YAML.
\usepackage{newunicodechar}
\newunicodechar{→}{\ensuremath{\rightarrow}}
% Mismo problema con los simbolos de marca y vinetas: el «✓» de «una palomita ✓»
% salia como cuadrito vacio en el PDF de todas las guias que lo usaban.
\usepackage{pifont}
\newunicodechar{✓}{\ding{51}}
\newunicodechar{✗}{\ding{55}}
\newunicodechar{✔}{\ding{52}}
\newunicodechar{•}{\textbullet}
\newunicodechar{≥}{\ensuremath{\geq}}
\newunicodechar{≤}{\ensuremath{\leq}}

\setmainfont{Helvetica Neue}
\setsansfont{Helvetica Neue}
\setlength{\parindent}{0pt}
\setlength{\parskip}{6pt}
% Sin particion de palabras: en 6.º-7.º una palabra cortada por guion al final
% de linea («Comuni-cacion») frena la lectura mucho mas que una linea corta.
\hyphenpenalty=10000
\exhyphenpenalty=10000
\pretolerance=2000
\tolerance=3000
\setlength{\emergencystretch}{3em}
\linespread{1.10}
\renewcommand{\arraystretch}{1.25}
\setlist{leftmargin=5mm,itemsep=1mm,topsep=1mm}
\newcolumntype{Y}{>{\RaggedRight\arraybackslash}X}

\definecolor{milcMagenta}{HTML}{E6007E}
\definecolor{milcVino}{HTML}{5A0038}
\definecolor{milcTurquesa}{HTML}{0093A5}
\definecolor{milcVerde}{HTML}{58A923}
\definecolor{milcMostaza}{HTML}{E5B400}
\definecolor{milcOcre}{HTML}{B86600}
\definecolor{milcNegro}{HTML}{241816}
\definecolor{milcGris}{HTML}{F7F7F4}
\definecolor{milcLinea}{HTML}{E8E2DD}
\definecolor{milcGradePrimary}{HTML}{0D9488}
\definecolor{milcGradeDark}{HTML}{115E59}
\definecolor{milcGradeSoft}{HTML}{CCFBF1}
\definecolor{milcGradeAccent}{HTML}{CFE7FF}
\definecolor{milcGradeText}{HTML}{FFFFFF}
\color{milcNegro}

\pagestyle{fancy}
\fancyhf{}
% Cabecera de dos lineas: identidad de la guia arriba y, a la derecha y en
% gris, la estacion en curso (\markright desde \milcestacion). Antes de la
% primera estacion la segunda linea va vacia; el \strut mantiene la altura
% para que la linea de arriba no baile entre paginas.
\lhead{\small Territorio interior · Educación socioemocional\\[-1pt]{\footnotesize\strut}}
\rhead{\small Grado 9° · Momento 5 · MILC v3\\[-1pt]{\footnotesize\strut\color{milcNegro!65}\rightmark}}
\cfoot{\small\thepage}
\renewcommand{\headrulewidth}{0pt}

% Sobre mostaza (#E5B400) y verde (#58A923) el texto va en milcNegro; sobre el
% resto de la paleta, en blanco. Se usa dentro de fonttitle/fontupper, donde un
% \color posterior manda sobre coltitle.
\newcommand{\milctintasobre}[1]{%
  \ifboolexpr{test{\ifstrequal{#1}{milcMostaza}} or test{\ifstrequal{#1}{milcVerde}}}%
    {\color{milcNegro}}{\color{white}}}

\newtcolorbox{titlebox}[1]{
  enhanced,colback=#1,colframe=#1,arc=7mm,boxrule=0pt,
  left=7mm,right=7mm,top=3mm,bottom=3mm,
  before skip=8pt,after skip=8pt,
  fontupper=\sffamily\bfseries\Large\color{white}
}

\newtcolorbox{softbox}[3]{
  enhanced,breakable,pad at break=3mm,
  colback=#2,colframe=#1,borderline west={4pt}{0pt}{#1},
  arc=2mm,boxrule=.4pt,left=4mm,right=4mm,top=3mm,bottom=3mm,
  before skip=6pt,after skip=8pt,title={#3},coltitle=white,
  colbacktitle=#1,fonttitle=\sffamily\bfseries\milctintasobre{#1},fontupper=\RaggedRight,
  attach boxed title to top left={xshift=3mm,yshift=-2mm},
  boxed title style={arc=2mm, boxrule=0pt}
}

\newtcolorbox{drawbox}[2]{
  enhanced,breakable,pad at break=4mm,
  colback=white,colframe=#1,arc=4mm,boxrule=1pt,
  left=5mm,right=5mm,top=4mm,bottom=4mm,before skip=8pt,after skip=8pt,
  title={#2},coltitle=white,colbacktitle=#1,fonttitle=\sffamily\bfseries\milctintasobre{#1},
  fontupper=\RaggedRight,
  attach boxed title to top left={xshift=4mm,yshift=-2mm},
  boxed title style={arc=3mm, boxrule=0pt}
}

\newtcolorbox{infoband}[1]{
  enhanced,breakable,pad at break=2mm,colback=#1!8,colframe=#1!50,arc=2mm,boxrule=.5pt,
  left=4mm,right=4mm,top=2mm,bottom=2mm,before skip=4pt,after skip=4pt,
  fontupper=\small\RaggedRight
}

% Puente narrativo entre secciones (contrato editorial Conéctate).
% Da continuidad: "Acabas de X. Ahora vas a Y porque Z."
% Visualmente: banda izquierda en color del grado, texto en itálica pequeña.
\newtcolorbox{puentebox}{
  enhanced,breakable,pad at break=2mm,colback=milcGradeSoft,colframe=milcGradeDark,
  borderline west={3pt}{0pt}{milcGradeDark},
  arc=0mm,boxrule=0pt,left=5mm,right=4mm,top=2.5mm,bottom=2.5mm,
  before skip=4pt,after skip=8pt,
  fontupper=\itshape\small\color{milcNegro}\RaggedRight
}
% La flecha va en modo matematico a proposito: Helvetica Neue NO tiene el glifo
% U+2192, asi que \textrightarrow se compilaba a NADA («Missing character: There
% is no → in font Helvetica Neue») y el puente salia sin flecha. En modo
% matematico la flecha viene de la fuente matematica y si se dibuja.
\newcommand{\puente}[1]{%
  \begin{puentebox}%
    \textcolor{milcGradeDark}{$\rightarrow$}\hspace{2mm}#1%
  \end{puentebox}%
}

\newcommand{\linead}[1]{\noindent\color{milcLinea}\rule{#1}{0.5pt}\color{milcNegro}}
\newcommand{\respuesta}[1]{\par\vspace{2mm}\foreach \n in {1,...,#1}{\noindent\color{milcLinea}\rule{\linewidth}{0.5pt}\par\vspace{5mm}}\color{milcNegro}}
\newcommand{\checkbox}{\textcolor{milcGradeDark}{\fbox{\phantom{x}}}\hspace{5pt}}

% ─── Listas aireadas para las guías socioemocionales ────────────────────
% Componer las verificaciones como «\checkbox texto\\» deja el texto que salta
% de línea debajo del cuadrito y apelmaza el bloque. Estos entornos alinean la
% continuación y airean los ítems. Los usa Territorio Interior; cualquier
% familia puede hacerlo.
\newlist{checklista}{itemize}{1}
\setlist[checklista]{
  label=\textcolor{milcGradeDark}{\fbox{\phantom{x}}},
  leftmargin=9mm, labelsep=3.2mm, itemsep=2.4mm,
  topsep=2.5mm, parsep=0pt, partopsep=0pt, before=\RaggedRight
}
\newlist{pasoslista}{enumerate}{1}
\setlist[pasoslista]{
  label=\textcolor{milcGradeDark}{\sffamily\bfseries\arabic*.},
  leftmargin=8mm, labelsep=3mm, itemsep=2.8mm,
  topsep=1mm, parsep=0pt, partopsep=0pt, before=\RaggedRight
}

\newcommand{\phasecard}[4]{
\begin{tcolorbox}[enhanced,colback=#3,colframe=#2,arc=3mm,boxrule=.5pt,height=3.05cm,valign=center,halign=center,left=1.5mm,right=1.5mm,top=2mm,bottom=2mm]
{\sffamily\bfseries\fontsize{22}{24}\selectfont\textcolor{#2}{#1}}\par
{\sffamily\bfseries\small #4}
\end{tcolorbox}
}

% ── Piezas del rediseno 2026 (legibilidad 6.º-11.º) ───────────────────
% Banda de estacion: reemplaza el titlebox macizo. Titulo a la izquierda,
% dato practico (tiempo/modalidad) a la derecha, en una sola linea.
% Nunca huerfana: pide 9 lineas libres (o salta de pagina rellenando con
% \vfill) y prohibe el salto justo despues de la banda. Nueve y no seis:
% la caja partible que sigue necesita ~80pt (titulo adosado, relleno y dos
% lineas) para arrancar en la misma pagina; con menos, tcolorbox la manda
% entera a la siguiente y la banda se queda sola. El penalty 10000 va
% en `after`, ANTES del espacio posterior: un `after skip` (aunque sea 0pt)
% es un \vskip tras la caja, y ese glue es un punto de corte legal que un
% \nopagebreak escrito despues de \end{tcolorbox} ya no cierra. Cada banda
% deja marcador PDF y actualiza la cabecera derecha.
\newcounter{milcestacion}
\newcommand{\milcestacion}[2]{%
\Needspace*{9\baselineskip}%
\stepcounter{milcestacion}%
\pdfbookmark[1]{#2}{milcestacion\themilcestacion}%
\markright{#2}%
\begin{tcolorbox}[enhanced,colback=#1,colframe=#1,arc=2mm,boxrule=0pt,
  left=5mm,right=5mm,top=2.6mm,bottom=2.6mm,before skip=11pt,
  after={\par\nopagebreak\vskip7pt}]
{\sffamily\bfseries\fontsize{15}{18}\selectfont\milctintasobre{#1}#2}
\end{tcolorbox}}

% Tarjeta de paso numerada: sustituye las tablas «Pilar / Aplicacion», que
% eran formato de consulta docente, no de estudiante. El numero va en un
% circulo sobre el borde para que la secuencia se lea de un vistazo.
\newcommand{\milcpaso}[3]{%
\Needspace{4\baselineskip}%
\begin{tcolorbox}[enhanced,colback=white,colframe=milcLinea,arc=1.5mm,boxrule=.9pt,
  left=14mm,right=4mm,top=3mm,bottom=3mm,before skip=4pt,after skip=4pt,
  fontupper=\RaggedRight,
  overlay={\node[anchor=north west,circle,fill=#2,text=white,inner sep=0pt,
    minimum size=7.4mm,font=\sffamily\bfseries\small]
    at ([xshift=3.6mm,yshift=-3.2mm]frame.north west) {#1};}]
#3
\end{tcolorbox}}

% Casilla marcable de tamano util con lapiz (la anterior era \fbox{\phantom{x}},
% de alto variable segun el texto que la seguia).
\newcommand{\milcmarca}{\framebox[3.8mm]{\rule{0pt}{3.4mm}}}

% Fila marcable de lista suelta (checks de la estacion 1).
\newcommand{\milccheckrow}[1]{%
\par\vspace{1.8mm}\noindent
\makebox[6.4mm][l]{\raisebox{-.55mm}{\textcolor{milcNegro}{\milcmarca}}}%
\parbox[t]{\dimexpr\linewidth-6.4mm\relax}{\RaggedRight #1}\par}

% Celda marcable dentro de la rubrica (tabla de autoevaluacion). Misma
% sangria francesa, pero pensada para vivir dentro de una columna X.
\newcommand{\milcopcion}[1]{%
\makebox[5.2mm][l]{\raisebox{-.55mm}{\textcolor{milcNegro}{\milcmarca}}}%
\parbox[t]{\dimexpr\linewidth-5.2mm\relax}{\RaggedRight #1}}

% ── Recursos visuales de la guía ──────────────────────────────────────
% Declarados en el bloque `recursos:` del YAML y emitidos por el builder.
% \guiaFigura   → imagen rasterizada o PDF (foto, carta celeste, montaje).
% \guiaDiagrama → fuente TikZ/circuitikz (.tex) incrustada como vector:
%                 es la vía para circuitos y esquemáticos editables.
% [#1] = texto alternativo (alt) · #2 = ruta relativa al .tex · #3 = pie
% El alt llega al PDF etiquetado (/Alt) y lo lee el lector de pantalla.
\newcommand{\guiaFigura}[3][]{%
\begin{center}
\includegraphics[width=0.88\linewidth,height=0.38\textheight,keepaspectratio,alt={#1}]{#2}\\[1.5mm]
{\footnotesize\itshape\textcolor{milcNegro!75}{#3}}
\end{center}
}

\newcommand{\guiaDiagrama}[2]{%
\begin{center}
\input{#1}\\[1.5mm]
{\footnotesize\itshape\textcolor{milcNegro!75}{#2}}
\end{center}
}

\begin{document}
\thispagestyle{empty}

% ── Portada ───────────────────────────────────────────────────────────
% Antes: un rectangulo de color a pagina completa. Se veia bien en pantalla,
% pero en la fotocopiadora del colegio se comia el toner y salia gris sucio.
% Ahora la banda ocupa el tercio superior y el resto va en blanco.
\begin{tikzpicture}[remember picture,overlay]
  \path[fill=milcGradePrimary]
    (current page.north west) rectangle ([yshift=-10.6cm]current page.north east);

  \node[anchor=north west,text=milcGradeText] at ([xshift=2.0cm,yshift=-1.55cm]current page.north west)
    {\sffamily\bfseries\fontsize{12}{14}\selectfont MOMENTO};
  \node[anchor=north west,text=milcGradeText,opacity=.85] at ([xshift=2.0cm,yshift=-2.35cm]current page.north west)
    {\sffamily\bfseries\fontsize{8.5}{10}\selectfont TERRITORIO INTERIOR · SOCIOEMOCIONAL};
  \node[anchor=north east,text=milcGradeText,opacity=.85,align=right] at ([xshift=-2.0cm,yshift=-1.55cm]current page.north east)
    {\sffamily\bfseries\fontsize{9}{12}\selectfont I.E. Sor María Juliana\\Cartago, Valle del Cauca};

  \node[anchor=west,text=milcGradeText] (gradonum) at ([xshift=1.85cm,yshift=-7.1cm]current page.north west)
    {\sffamily\bfseries\fontsize{112}{112}\selectfont 5};
  % El circulito de grado (°) solo aplica a las guias de grado («11°»); en
  % Bebras y Semillero el numero grande es un momento/modulo, no un grado.
  % Se posiciona relativo al nodo `gradonum`, no en coordenadas absolutas.
  

  \node[anchor=west,text=milcGradeText,text width=11.4cm,align=left] at ([xshift=1.5cm,yshift=0.5cm]gradonum.east)
    {
     {\sffamily\bfseries\fontsize{9}{11}\selectfont\MakeUppercase{Territorio interior · Socioemocional}}\\[3mm]
     {\sffamily\bfseries\fontsize{21}{25}\selectfont\setlength{\baselineskip}{27pt}La broma\\[4pt]que no es broma\\[4pt]el color tuvo precio de lista}\\[3.5mm]
     {\sffamily\bfseries\fontsize{15}{18}\selectfont Grado 9° · Momento 5}
    };
\end{tikzpicture}

\vspace*{9.4cm}
\vfill

{\sffamily\bfseries\fontsize{8}{10}\selectfont\textcolor{milcGradeDark}{LO QUE TE LLEVAS HOY}}

\begin{tcolorbox}[enhanced,colback=white,colframe=milcNegro,arc=2mm,boxrule=1.6pt,
  left=5mm,right=5mm,top=4mm,bottom=4mm,before skip=3pt,after skip=12pt,
  fontupper=\RaggedRight\fontsize{11}{15.5}\selectfont]
Tu registro de una semana: las frases sobre color, pelo, origen o plata que oíste —con quién las dijo sin ninguna mala intención—, clasificadas en los tres tipos, y las dos respuestas que preparaste: una para el que la dijo y otra para el que la recibió.
\end{tcolorbox}

\begin{tcolorbox}[enhanced,colback=milcGris,colframe=milcGris,arc=2mm,boxrule=0pt,
  left=5mm,right=5mm,top=3.5mm,bottom=3.5mm,after skip=12pt]
{\sffamily\bfseries\fontsize{8}{10}\selectfont\textcolor{milcNegro!55}{ESTA GUÍA ES DE}}\\[3mm]
\begin{tabularx}{\linewidth}{X p{3.2cm} p{3.2cm}}
\linead{\linewidth} & \linead{3.0cm} & \linead{3.0cm}\\[-1mm]
{\fontsize{8.5}{10}\selectfont\textcolor{milcNegro!55}{Nombre completo}} &
{\fontsize{8.5}{10}\selectfont\textcolor{milcNegro!55}{Grupo}} &
{\fontsize{8.5}{10}\selectfont\textcolor{milcNegro!55}{Fecha}}\\
\end{tabularx}
\end{tcolorbox}

\vfill

\noindent\textcolor{milcLinea}{\rule{\linewidth}{1pt}}\\[2mm]
{\sffamily\bfseries\fontsize{9.5}{12}\selectfont Autor: Dr. Álvaro Cárdenas Orozco}\\[1mm]
{\sffamily\fontsize{8.5}{11}\selectfont\textcolor{milcNegro!55}{Metodología MILC · Apertura ancestral · Racismo y clasismo cotidianos · Triángulo Dussel-Estoicismo-Floridi}}

\newpage

\section*{Apertura: La broma que no es broma: cuando el color y la plata tuvieron precio de lista}

\begin{softbox}{milcGradeDark}{milcGradeSoft}{Saber ancestral}
El \textbf{10 de febrero de 1795}, la Corona española expidió para América una \textbf{Real Cédula de Gracias al Sacar}. \emph{No era un discurso ni una doctrina:} \textbf{era un arancel} ---una lista de precios--- donde se decía \textbf{cuánto había que pagar} para obtener \textbf{setenta y una} clases de gracias, dispensas y trámites. \textbf{Y entre esas setenta y una opciones estaba \emph{la blancura}.} \emph{Los precios están documentados:} la dispensa de la \textbf{calidad de pardo} costaba \textbf{100 reales}; la de \textbf{quinterón}, \textbf{800}; y el derecho a usar el título de \textbf{«don»}, \textbf{1.000}. \textbf{Pardos y mulatos podían, literalmente, \emph{comprar} su blancura} ---y con ella el acceso a ciertas instituciones educativas y a cargos públicos---. \emph{La élite criolla protestó, porque eso les daba a los pardos derechos que hasta entonces solo tenían los blancos.} \textbf{Guarda esto, porque es la idea que sostiene toda la guía:} \emph{el color no era una descripción de la piel ---era \textbf{una posición}, y esa posición tenía un precio y unas consecuencias muy concretas: qué podías estudiar, qué cargo podías ocupar, cómo te decían.} \textbf{Nadie tiene ya esa lista. Pero la posición no desapareció: se volvió invisible.}
\end{softbox}

\begin{softbox}{milcTurquesa}{milcGris}{Saber contemporáneo}
¿Y cómo se sostiene una jerarquía cuando ya nadie la defiende en voz alta? \textbf{Con cosas pequeñas y constantes.} El psicólogo \textbf{Derald Wing Sue} y su equipo estudiaron lo que llamaron \textbf{microagresiones raciales}: \emph{indignidades \textbf{breves y cotidianas} ---verbales, conductuales o del entorno---, \textbf{intencionales o no}, que comunican desprecio o insulto hacia personas racializadas} (Sue y otros, 2007). \textbf{Y las clasificaron en tres tipos.} El \textbf{microasalto} es el más parecido al racismo de antes: explícito, deliberado, un insulto. El \textbf{microinsulto} es una grosería disfrazada, muchas veces con forma de cumplido ---«qué bien hablas», «eres bonita \emph{para ser} morena»---. Y la \textbf{microinvalidación} borra la experiencia del otro: «eso no fue racismo», «estás exagerando», «aquí no somos racistas». \textbf{Ahora el hallazgo que hay que subrayar, porque es el que cambia la conversación:} \emph{quienes cometen microagresiones \textbf{con frecuencia no son conscientes} de estarlo haciendo}. \textbf{Es decir: la ausencia de mala intención no es una excepción rara ---es \emph{lo normal}---.} \emph{Y por eso «yo no lo dije con mala intención» describe correctamente el caso y no lo resuelve.}
\end{softbox}

\begin{softbox}{milcMostaza}{milcGris}{Pregunta-puente}
\textit{En 1795 la blancura estaba en \textbf{una lista, con precio y con firma}. Hoy no hay lista. \textbf{¿Qué frases sobre el color, el pelo, el acento o la plata circulan en tu curso} \ldots \textbf{dichas por gente que jura no tener nada contra nadie}?}
\end{softbox}

\begin{softbox}{milcVerde}{milcGris}{Saber y saber hacer}
Al terminar podrás: (1) \textbf{entender} que las jerarquías de color y clase en Colombia fueron \textbf{norma escrita}, con precios documentados, y no un malentendido de las personas; (2) \textbf{clasificar} frases reales en microasalto, microinsulto y microinvalidación; (3) \textbf{explicar} por qué la falta de intención no cancela el daño y por qué el «elogio» puede ser el caso más difícil; (4) \textbf{responder} de forma corta a quien lo dice y \textbf{acompañar} a quien lo recibe.
\end{softbox}



\small
\begin{tabularx}{\textwidth}{>{\bfseries}p{3.0cm}Y}
\rowcolor{milcGradePrimary!12}Autor & Dr. Álvaro Cárdenas Orozco\\
DBA & Comprende los fenómenos que afectan a su territorio, regula sus emociones ante ellos y acompaña a otros con empatía (transversal 6°--11°, articulado con la política «Escuela, territorio de vida» del MEN, Resolución 006519 de 2025, y la Ley 1620 de 2013)\\
Referentes & Primera ayuda psicológica (OMS, 2011/2012) · Directrices IASC (2007) · USGS y Servicio Geológico Colombiano · Pueblo nasa y la minga (CRIC) · Dussel · Estoicismo · Floridi\\
Producto final & Tu registro de una semana: las frases sobre color, pelo, origen o plata que oíste —con quién las dijo sin ninguna mala intención—, clasificadas en los tres tipos, y las dos respuestas que preparaste: una para el que la dijo y otra para el que la recibió.\\
\end{tabularx}
\normalsize

\puente{En el momento 4 del grado 7 viste la palabra que hiere. \textbf{Hoy es la misma herramienta apuntando a algo que nadie escoge:} \emph{el color, el pelo, el acento, de dónde viene su familia, cuánto hay en la casa}.}

\milcestacion{milcGradeDark}{Ruta MILC: así vamos a aprender}
\begin{tabularx}{\textwidth}{>{\centering\arraybackslash}X>{\centering\arraybackslash}X>{\centering\arraybackslash}X>{\centering\arraybackslash}X}
\phasecard{1}{milcVerde}{milcGris}{Escucha\\[-1mm]\small Observo, escucho y pregunto.} &
\phasecard{2}{milcTurquesa}{milcGris}{Sistematización\\[-1mm]\small Nombro lo que no se nota.} &
\phasecard{3}{milcMagenta}{milcGris}{Praxis\\[-1mm]\small Construyo y aplico.} &
\phasecard{4}{milcVino}{milcGris}{Evaluación\\[-1mm]\small Reflexión liberadora.}
\end{tabularx}


\puente{Primero, sin acusar a nadie: qué se dice, y qué se dice \emph{con cariño}, que es lo más difícil de ver.}

\milcestacion{milcVerde}{Estación 1 de 3 · Escucha}

\begin{softbox}{milcVerde}{milcGris}{Escena para escuchar}
\textbf{Actividad 1 · IDENTIFICA --- Lo que se dice.} \textbf{Sin nombres: usa iniciales.} \textbf{Primera parte.} Escribe \textbf{las frases que circulan en tu colegio} sobre el color de piel, el pelo, los rasgos, el acento, de dónde viene la familia de alguien o cuánta plata hay en su casa. \emph{Incluye apodos.} \textbf{Escríbelas tal cual, aunque incomode verlas escritas ---y esa incomodidad ya es información---.} \textbf{Segunda parte, la difícil.} Marca cuáles se dicen \textbf{«de cariño»} o \textbf{como elogio}: «qué bonito tu pelo, ¿te lo puedo tocar?», «hablas muy bien para ser de allá», «tú eres el negro pero eres bien». \textbf{Tercera parte.} ¿Qué pasa cuando alguien dice que eso le molestó? \emph{Anota la respuesta textual: «era chiste», «qué delicado», «no sea exagerado», «aquí nadie es racista».} \textbf{Y la última, para ti solo:} ¿has dicho alguna de estas frases? \emph{¿Qué querías decir en ese momento?}
\end{softbox}

{\sffamily\bfseries\fontsize{8}{10}\selectfont\textcolor{milcNegro!55}{MARCA CUANDO LO HAYAS HECHO}}
\milccheckrow{Escribiste las frases tal cual, incluidos los apodos.}
\milccheckrow{Marcaste las que se dicen \textbf{de cariño o como elogio}.}
\milccheckrow{Anotaste textual qué responden cuando alguien dice que le molestó.}

\begin{infoband}{milcVerde}


\textbf{En tu cuaderno (Actividad 1 · IDENTIFICA).} Título: «Actividad 1. Lo que se dice». \textbf{Formato:} las frases, cuáles son «de cariño», qué responden cuando alguien reclama, y tu propia participación. \textbf{Extensión:} media página. \textbf{Sabes que terminaste cuando} tienes \textbf{al menos una frase que se dice como elogio}. Esta página es tuya: \textbf{nadie más tiene que leerla si tú no quieres}.
\end{infoband}

\puente{Ya lo tienes. Ahora, que esto fue norma escrita con precio, y los tres tipos en que se presenta hoy.}

\milcestacion{milcTurquesa}{Fase 2 · Sistematización: entender la microagresión}

\begin{softbox}{milcTurquesa}{milcGris}{Cuatro claves sobre el racismo cotidiano}
Cuatro claves para dejar de discutir si alguien «es racista» y empezar a mirar lo que pasa realmente.
\end{softbox}

\milcpaso{1}{milcTurquesa}{\textbf{Tuvo precio de lista, y eso lo explica casi todo.} La Real Cédula de \textbf{Gracias al Sacar} (1795) puso la blancura entre \textbf{setenta y una} opciones de compra: \textbf{100 reales} la dispensa de pardo, \textbf{800} la de quinterón, \textbf{1.000} el título de «don». \emph{Lo que se compraba no era un color de piel ---eso no cambia--- sino \textbf{el permiso de entrar}:} a ciertos estudios, a ciertos cargos, a cierto trato. \textbf{Y la élite criolla protestó, lo cual demuestra que el asunto era real y no simbólico.} \emph{Conclusión que hay que llevarse:} \textbf{el racismo en América no empezó como una antipatía personal ---empezó como una \textbf{norma que repartía derechos}---.} \emph{Por eso no se arregla con que la gente sea más amable, y por eso lo que hoy queda no son opiniones sueltas: son restos de un reparto.}}
\milcpaso{2}{milcTurquesa}{\textbf{Los tres tipos, y el peligroso es el segundo.} \textbf{Microasalto:} explícito y deliberado ---el insulto, el apodo racial, la exclusión abierta---. \textbf{Microinsulto:} una grosería con forma de comentario amable, y muchas veces \emph{de elogio} ---«qué bien hablas», «bonita \emph{para ser} morena», «no pareces de allá»---. \textbf{Microinvalidación:} borrar la experiencia del otro ---«eso no fue racismo», «estás exagerando», «aquí no somos así»--- (Sue y otros, 2007). \emph{Fíjate en la trampa del segundo:} \textbf{todo elogio que lleve un «para ser» adentro está diciendo, en el fondo, «lo esperado de los tuyos era menos»}. \textbf{Y fíjate en el tercero, porque es el que más duele:} \emph{la microinvalidación no le hace daño a alguien por lo que dice, sino porque lo deja \textbf{sin poder nombrar lo que le pasó}}.}
\milcpaso{3}{milcTurquesa}{\textbf{Casi nadie sabe que lo está haciendo.} Sue y su equipo lo señalan con claridad: \emph{quienes cometen microagresiones \textbf{con frecuencia no son conscientes} de hacerlo}. \textbf{Eso hay que tomarlo en serio en las dos direcciones.} \emph{Hacia el que lo dice:} \textbf{no es un monstruo, y tratarlo como tal garantiza que se defienda y no cambie nada}. \emph{Hacia el que lo escucha:} \textbf{que no haya intención no significa que no haya daño}, y la pregunta útil no es «¿fue racista esa persona?» sino «\emph{¿qué le hizo esa frase a quien la recibió?}». \textbf{Es la misma lógica del momento 4 del grado 7:} \emph{«era chiste» describe la intención y no toca la consecuencia.}}
\milcpaso{4}{milcTurquesa}{\textbf{El costo lo paga siempre el que recibe.} \emph{Una frase suelta parece nada.} \textbf{Pero la persona que la recibe no recibe una: recibe la de hoy, la de ayer y la de la semana pasada,} y cada vez tiene que decidir en dos segundos si dice algo ---y quedar de amargada--- o se la aguanta. \textbf{Y si dice algo, viene la microinvalidación,} \emph{así que muchas veces el cálculo racional es callarse}. \textbf{Por eso el trabajo no puede recaer en quien lo recibe.} \emph{Aquí vale exactamente lo del momento 3 del grado 8:} \textbf{quien tiene que hacer algo es el que está mirando}, porque a él no le cuesta casi nada y no está cansado de esto.}

\begin{drawbox}{milcTurquesa}{Cómo se responde, según quién seas en la escena}
\begin{pasoslista}
\item \textbf{Si lo dijiste tú:} «tienes razón, eso sonó mal» y sigue. \emph{Sin explicar tus intenciones ---a nadie le sirven---.}
\item \textbf{Si lo dijo otro y estabas ahí:} corto y sin sermón: «uy, con eso no» · «esa no».
\item \textbf{Si te lo dijeron a ti:} \emph{no estás obligado a educar a nadie}. Si quieres: «no me gusta que me digan así».
\item \textbf{Nunca:} «no era con mala intención» como respuesta. \emph{Es cierto y no viene al caso.}
\item \textbf{Después:} búscalo en privado. \emph{Corregir delante de todos casi siempre produce defensa, no cambio.}
\end{pasoslista}
\end{drawbox}

\begin{softbox}{milcMostaza}{milcGris}{Errores comunes y su corrección}
\textbf{«No fue con mala intención»:} describe la intención y no toca el daño. \textbf{«Aquí nadie es racista»:} es una microinvalidación de manual. \textbf{Discutir si alguien \emph{es} racista:} la pregunta útil es qué hizo la frase. \textbf{Los elogios con «para ser»:} dicen que de los suyos se esperaba menos. \textbf{Pedirle a quien lo recibe que explique:} ya está cansado; no es su trabajo. \textbf{Tocarle el pelo a alguien:} preguntarlo no lo vuelve aceptable. \textbf{El sermón público:} produce defensa, no cambio.
\end{softbox}

\begin{infoband}{milcTurquesa}


\textbf{En tu cuaderno (Actividad 2 · EXPLICA).} Título: «Actividad 2. Los tres tipos». \textbf{Formato:} clasifica \textbf{seis} frases de tu Actividad 1 en microasalto, microinsulto y microinvalidación, y explica por qué cada una. \textbf{Extensión:} media página. \textbf{Sabes que terminaste cuando} puedes explicar por qué \textbf{un elogio con «para ser» adentro es un insulto}.
\end{infoband}

\puente{Ya sabes clasificarlo. Es hora de preparar dos respuestas: una para el que lo dijo sin darse cuenta y otra para el que lo recibió.}

\milcestacion{milcMagenta}{Fase 3 · Praxis: responder sin escándalo}

\begin{softbox}{milcMagenta}{milcGris}{Cómo se responde a algo que «no era con mala intención»}
Una semana de registro y dos frases preparadas. Eso es todo, y es más de lo que la mayoría hace en años.
\end{softbox}

\milcpaso{1}{milcMagenta}{\textbf{El registro de la semana (7 días, 2 min al día).} Anota \textbf{cada frase} que oigas sobre color, pelo, rasgos, acento, origen o plata ---sin nombres---, y al lado \textbf{si quien la dijo parecía tener mala intención}. \emph{Al final de la semana, cuenta cuántas fueron sin ninguna mala intención.} \textbf{Ese número es el argumento entero de la guía.}}
\milcpaso{2}{milcMagenta}{\textbf{Clasificar los tres tipos (8 min).} Pasa tu registro a las tres categorías. \emph{Presta atención a cuál es la más numerosa:} \textbf{casi siempre gana el microinsulto}, y casi siempre viene disfrazado de comentario simpático.}
\milcpaso{3}{milcMagenta}{\textbf{La respuesta corta (10 min, en tríos).} Escribe y ensaya \textbf{dos frases}: una para \textbf{cortar} en el momento ---menos de cinco segundos, sin humillar a quien lo dijo, como en el momento 4 del grado 7--- y otra para \textbf{acompañar} después, en privado, a quien lo recibió. \emph{Ensáyenlas en voz alta con casos inventados.} \textbf{Y observen algo:} \emph{la que cuesta no es la primera. La que cuesta es sostener el silencio incómodo de los dos segundos que siguen.}}
\milcpaso{4}{milcMagenta}{\textbf{El elogio que ofende (8 min).} Tomen \textbf{tres «elogios»} del registro y reescríbanlos \textbf{quitándoles la comparación}: \emph{«bonita para ser morena» → «bonita»; «hablas muy bien para ser de allá» → «hablas muy bien»}. \emph{Léanlos en las dos versiones, en voz alta.} \textbf{La diferencia se oye,} \emph{y ese ejercicio de treinta segundos hace más que cualquier explicación}.}

\begin{drawbox}{milcMagenta}{Antes de cerrar tu registro}
\begin{checklista}
\item Registraste \textbf{siete días} y anotaste si parecía haber mala intención.
\item Contaste cuántas frases fueron dichas \textbf{sin ninguna mala intención}.
\item Clasificaste seis frases en los tres tipos, con la razón.
\item Tu frase para cortar cabe en cinco segundos y \textbf{no humilla} a quien lo dijo.
\item Tu frase para acompañar es en privado y \textbf{no le pide nada} a quien lo recibió.
\item Reescribiste tres «elogios» quitándoles la comparación y los leíste en voz alta.
\end{checklista}
\end{drawbox}

\begin{softbox}{milcMagenta}{milcGris}{Plantilla de trabajo}
\textbf{Para cortar.} \emph{«Uy, con eso no» · «esa no» · «no le veo la gracia».} \\[1mm]
\textbf{Para acompañar, en privado.} \emph{«Oí lo que te dijeron. Me pareció una mala jugada.»} ---sin consejo, sin pedirle nada---. \\[1mm]
\textbf{Si el que lo dijo fui yo.} \emph{«Tienes razón, eso sonó mal.»} ---y seguir; sin explicar intenciones---.
\end{softbox}

\begin{infoband}{milcMagenta}


\textbf{En tu cuaderno (Actividad 3 · CREA).} Título: «Actividad 3. Mi registro de la semana». \textbf{Formato:} el registro de siete días con el conteo + la clasificación en tres tipos + las dos frases + los tres elogios reescritos. \textbf{Extensión:} una página. \textbf{Sabes que terminaste cuando} tienes \textbf{el número de frases dichas sin mala intención}.
\end{infoband}

\puente{El producto es un registro de una semana. \emph{Un registro convence de una manera que ninguna clase logra: porque lo escribió uno mismo.}}

\milcestacion{milcMostaza}{Producto de la sesión}
\textbf{Registro de una semana: frases clasificadas en los tres tipos y dos respuestas preparadas}.
\begin{softbox}{milcMostaza}{milcGris}{Criterios de calidad}
(1) El registro cubre siete días e indica si parecía haber mala intención en cada caso. (2) La clasificación en microasalto, microinsulto y microinvalidación está justificada frase por frase. (3) Incluye frases dichas «de cariño» o como elogio, no solo insultos evidentes. (4) La respuesta para cortar es breve y no humilla a quien habló; la de acompañar no pide nada a quien recibió. (5) Se reescriben elogios quitándoles la comparación, mostrando qué cambia.
\end{softbox}

\puente{Antes de cerrar, revisa lo que casi siempre falta: si en tu registro no hay ninguna frase dicha «de cariño», lo hiciste incompleto.}

\milcestacion{milcVino}{Evaluación liberadora · cinco dimensiones}

\begin{softbox}{milcVino}{milcGris}{Sentido de la evaluación}
Evaluar no es solo revisar una nota. Es reconocer qué aprendí, qué cambió en mí, qué emociones manejé, cómo me relaciono con la ciudad y la comunidad, y qué le dejaré a quien viene detrás.
\end{softbox}

\textbf{Mi reflexión liberadora en cinco dimensiones.} Completa cada frase iniciada:

\small
\begin{tabularx}{\textwidth}{>{\bfseries\RaggedRight\arraybackslash}p{3.4cm}Y>{\RaggedRight\arraybackslash}p{4.6cm}}
\rowcolor{milcVino}\color{white}Dimensión & \color{white}\bfseries Mi respuesta (frame starter) & \color{white}\bfseries Algo que descubrí\\
1. Personal & Lo que más cambió en mí fue \linead{6.5cm} & \linead{4.2cm}\\
2. Emocional & La emoción más fuerte fue \linead{2.5cm} y la regulé \linead{3cm} & \linead{4.2cm}\\
3. Ciudadana & En democracia esto importa porque \linead{6cm} & \linead{4.2cm}\\
4. Local (Cartago) & En mi barrio sería útil para \linead{6.5cm} & \linead{4.2cm}\\
5. Intergeneracional & A mi abuela/abuelo le diría \linead{6.5cm} & \linead{4.2cm}\\
\end{tabularx}
\normalsize

\begin{infoband}{milcVino}
\textbf{Para la próxima sustentación.} Vuelve a esta tabla en seis meses. Verás que las respuestas cambian — no porque lo aprendido cambie, sino porque tú cambias. La evaluación liberadora es una conversación contigo a través del tiempo.
\end{infoband}


\puente{Cerramos con tres voces: el otro que no cabe en una categoría, la diferencia entre la cosa y la opinión, y los datos pequeños con que se arma un juicio.}

\milcestacion{milcGradeDark}{Triángulo de pensamiento · cierre filosófico}

\begin{softbox}{milcVino}{milcGris}{Dussel · lente del nosotros}
\textit{«Analéctico quiere indicar el hecho real humano por el que todo hombre, todo grupo o pueblo se sitúa siempre más allá (aná-) del horizonte de la totalidad.»}

{\footnotesize\itshape\color{milcNegro!70}--- Enrique Dussel · Filosofía de la liberación (1977), §2.4}

La Gracias al Sacar es \textbf{una totalidad en su forma más pura}: un sistema que decidió qué era una persona ---pardo, quinterón, blanco--- \textbf{y le puso precio a cambiar de casilla}. \emph{Y fíjate en la trampa perfecta que eso arma:} \textbf{permitía la movilidad individual sin tocar la jerarquía}. \emph{Uno podía comprarse la blancura; la lista seguía intacta.} \textbf{Y algo parecido pasa hoy con los elogios:} \emph{«tú eres distinto», «tú sí eres bien» sacan a una persona de la categoría \textbf{y dejan la categoría en su sitio}} ---es decir, no son un desmentido del prejuicio: son su confirmación con una excepción---. \textbf{Lo analéctico dice otra cosa:} \emph{nadie cabe en la casilla, y por eso la casilla es lo que sobra.}

\textbf{¿He dicho «tú eres distinto» como si fuera un cumplido? ¿Distinto de qué, exactamente?}
\end{softbox}
\linead{\linewidth}\\[1mm]
\linead{\linewidth}

\begin{softbox}{milcMostaza}{milcGris}{Estoicismo · Epicteto · Enquiridión, 5 (c. 125 d.C.) · cuidado interior}
\textit{«No son las cosas las que perturban a los hombres, sino las opiniones sobre las cosas.»}

{\footnotesize\itshape\color{milcNegro!70}--- Epicteto · Enquiridión, 5 (c. 125 d.C.)}

\textbf{Cuidado con esta frase, porque se usa mal todo el tiempo.} \emph{Mal uso:} decirle a alguien que si le duele un comentario racista es por su opinión sobre el comentario ---eso es una microinvalidación con toga de filósofo---. \textbf{Uso correcto, y es potente:} \emph{lo que llamamos «color» no es una cosa ---es \textbf{una opinión sobre una cosa}---}. \textbf{La piel es un hecho; que la piel signifique algo sobre quién eres es una opinión}, y en 1795 esa opinión tenía tarifa. \emph{Y como es una opinión, se puede examinar y se puede cambiar} ---cosa que no pasaría si fuera un hecho de la naturaleza---. \textbf{Ahí está la esperanza de esta guía:} \emph{lo que se construyó se puede desarmar.}

\textbf{¿Qué cosas creo saber de alguien apenas con verlo? ¿De dónde saqué eso?}
\end{softbox}
\linead{\linewidth}\\[1mm]
\linead{\linewidth}

\begin{softbox}{milcTurquesa}{milcGris}{Floridi · lente de la infoesfera}
\textit{«Los pequeños patrones solo pueden ser significativos si se agregan correctamente, se comparan y se procesan a tiempo.»}

{\footnotesize\itshape\color{milcNegro!70}--- Luciano Floridi · Big data and their epistemological challenge (2012)}

Floridi describe, sin proponérselo, exactamente cómo funciona una microagresión: \textbf{es un \emph{patrón pequeño}}. \emph{Aislado no significa nada ---por eso quien lo dice puede decir con sinceridad que no fue nada---.} \textbf{El significado aparece en la \emph{agregación}: veinte frases pequeñas, de veinte personas distintas, a lo largo de un año, sobre la misma persona.} \emph{Y aquí está la asimetría que hay que entender:} \textbf{el que la dice ve un dato; el que la recibe ve la serie completa}. \emph{Por eso las dos personas pueden estar diciendo la verdad al mismo tiempo cuando una dice «no fue nada» y la otra dice «estoy harta».} \textbf{Y por eso el registro de una semana funciona:} \emph{te deja ver la serie que normalmente solo ve uno.}

\textbf{Si me dijeran una de esas frases una vez al día durante un año, ¿seguiría pareciéndome poca cosa?}
\end{softbox}
\linead{\linewidth}\\[1mm]
\linead{\linewidth}

\begin{infoband}{milcNegro}
\textbf{Nota para el docente.} Las tres son citas textuales y llevan su referencia debajo. Puedes remitir a la obra en clase.
\end{infoband}

\milcestacion{milcVino}{Mi autoevaluación · marca una casilla por fila}

{\small
\setlength{\extrarowheight}{2.2pt}
\begin{tabularx}{\textwidth}{>{\RaggedRight\arraybackslash\bfseries}p{3.0cm}YYY}
\rowcolor{milcVino}\color{white}\bfseries Criterio & \color{white}\bfseries Lo logré & \color{white}\bfseries En proceso & \color{white}\bfseries Necesito apoyo\\[.6mm]
Comprensión del tema & \milcopcion{Explico las ideas con mis palabras.} & \milcopcion{Reconozco las ideas, falta precisión.} & \milcopcion{Necesito repasar lo básico.}\\[.8mm]
\rowcolor{milcGris}Calidad de la respuesta & \milcopcion{Distingo los tres tipos y tengo una respuesta corta que no humilla a quien la dijo.} & \milcopcion{Reconozco el insulto directo, pero no las frases que parecen elogios.} & \milcopcion{Sigo pensando que si no hubo mala intención no hubo problema.}\\[.8mm]
Aplicación y producto & \milcopcion{Mi producto cumple los criterios.} & \milcopcion{Está armado, con ajustes pendientes.} & \milcopcion{Aún está en construcción.}\\[.8mm]
\rowcolor{milcGris}Reflexión 5 dimensiones & \milcopcion{Respondí cada una con honestidad.} & \milcopcion{Respondí algunas con detalle.} & \milcopcion{Mis respuestas son muy generales.}\\[.8mm]
Triángulo de pensamiento & \milcopcion{Conecté las 3 voces con mi trabajo.} & \milcopcion{Conecté al menos una.} & \milcopcion{No las relaciono con mi proceso.}\\[.8mm]
\end{tabularx}}

\vspace{3mm}
\textbf{Hoy entendí que:}\\[1mm]
\linead{\linewidth}\\[1mm]
\linead{\linewidth}

\textbf{Mi compromiso para la próxima sesión es:}\\[1mm]
\linead{\linewidth}\\[1mm]
\linead{\linewidth}

\begin{softbox}{milcMostaza}{milcGris}{Cita de despedida · Marco Aurelio}
\textit{``El obstáculo en el camino se vuelve el camino. Lo que estorba al camino se vuelve camino.''}\\[1mm]
Cada error de hoy, bien observado, se convierte en pieza del camino que pisarás mañana.
\end{softbox}

\small
\begin{tabularx}{\textwidth}{>{\bfseries}p{3.6cm}Y}
\rowcolor{milcGradeDark!12}Nombre completo: & \linead{10cm}\\
Fecha: & \linead{4cm} \quad Curso/grupo: \linead{4cm}\\
Mi firma: & \linead{9cm}\\
\end{tabularx}
\normalsize

\begin{infoband}{milcGradeDark}
\textbf{Para la próxima sesión.} Dos cosas. \textbf{Primero, trae el registro con el conteo} ---sobre todo el número de frases dichas \textbf{sin mala intención}--- y usa \textbf{una vez} tu frase para cortar. \emph{Anota qué pasó en los dos segundos siguientes, que es lo único difícil de todo esto.} \textbf{Segundo, pregunta en tu casa por lo que se decía:} averigua con alguien mayor \textbf{qué palabras se usaban en su época para hablar del color de la piel o del origen de la gente}, \textbf{si se decían como insulto o como algo normal}, y \textbf{si alguna vez alguien de la familia tuvo dificultades por eso} ---para entrar a un lugar, para un trabajo, para un noviazgo---. \textbf{Trae:} tu conteo y lo que te contaron. \emph{Prepárate para dos respuestas frecuentes y las dos enseñan: unos van a decir que eso «se decía sin maldad», y otros van a contar algo muy concreto que les pasó. Las dos cosas son ciertas a la vez, y esa simultaneidad es justamente el tema de esta guía.}
\end{infoband}

\end{document}
